% ---------------------------------------------------------------------------
%  I. Introduction
%  Floats:  figures/fig-robot.tex   (Fig. 1, robot photograph)
%  Holds:   the three structural obstacles, the research question,
%           H1-H4 as statements (verdicts belong in section 06),
%           the three contributions, the ethics pointer.
% ---------------------------------------------------------------------------

\section{Introduction}
\label{sec:intro}

A laboratory that buys a commercial service robot buys a sealed appliance. The vendor ships an
Android application, a cloud account, and a support contract; it does not ship the message
schemas that would let the owner send the robot anywhere the application does not already go.
For research groups this is a hard limit: the platform in Fig.~\ref{fig:robot} navigates,
speaks and greets visitors perfectly well, but only in the ways its manufacturer anticipated.

% Fig. 1 — the study platform, photograph. Used by sections/01-introduction.
% Source image: assets/robot_kiosk.jpg (institution logo, on-screen name and
% QR code are blurred; regenerate with tools/prepare_assets.py if re-cropped).
\begin{figure}[!t]
\centering
\includegraphics[width=0.62\columnwidth]{robot_kiosk}
\caption{\Platform{} running our reconstructed visitor interface on its own touchscreen.
The head is an Android computer; the mobile base underneath is a separate Linux/\ros{} machine.
Neither exposes a documented interface to its owner. Identifying marks are blurred.}
\label{fig:robot}
\end{figure}


\textbf{Why this is hard ?} Three obstacles compound, and none of them is peculiar to this
vendor. First, the robot is not one computer but two: an Android head and a \ros{} chassis,
joined by an undocumented internal link. No single point of observation covers the whole command
surface, and each one is blind in a different place (Section~\ref{sec:method}). Second, there is
no white-box route. Shells on the chassis are closed, and repeated authentication attempts lock
the account, so we cannot simply read the running system. Third, and least obvious, the software
layer that accepts a message is not the layer that turns the wheels. A command can be accepted,
acknowledged and then quietly dropped. Treat the acknowledgment as proof and the result is a
protocol description that is confidently wrong. This is the reason naive fuzzing of the exposed
interfaces gets nowhere on such a platform.

\textbf{Research question.} \emph{Can a non-destructive methodology recover the command surface
of a closed Android--\ros{} robot without documentation, and what evidence is required before a
recovered command may be considered validated?}

We formulate four hypotheses, stated here and resolved in Section~\ref{sec:eval}:

\begin{itemize}[leftmargin=1.2em]
\item \textbf{H1 (application-as-specification).} Static analysis of the vendor's own Android
application yields the wire protocol the chassis actually speaks.
\item \textbf{H2 (broker-as-command-channel).} The exposed MQTT broker can issue motion
commands, as the WebSocket bridge does.
\item \textbf{H3 (map-artifact reuse).} Reusing the vendor's own map annotations is more
reliable than re-deriving goal coordinates offline.
\item \textbf{H4 (transport $\neq$ execution).} A transport acknowledgment does not imply that
the chassis executed the command.
\end{itemize}

\textbf{Contributions.}
\begin{enumerate}[label=\textbf{C\arabic*.},leftmargin=1.5em]
\item A reproducible seven-phase reverse-engineering methodology for closed Android--\ros{}
robots, organised around falsifiable hypotheses and a triangulation rule, together with the
effort it required in practice (Sections~\ref{sec:method} and~\ref{sec:eval}).
\item A recovered command surface for a commercial reception platform and an edge-first
architecture built on it, including a speech channel obtained through Android inter-process
communication after every protocol-level channel had been eliminated
(Section~\ref{sec:platform}).
\item Field evidence that execution must be verified against state telemetry rather than
transport acknowledgments, with the resulting design rule, quantified indicators, and an
analysis of which findings should transfer to other closed platforms
(Sections~\ref{sec:eval} and~\ref{sec:discussion}).
\end{enumerate}

All work reported here was carried out on hardware owned by our institution, for
interoperability purposes, without redistributing vendor code or circumventing any licensing
mechanism; Section~\ref{sec:discussion} details the ethical and legal framing.
